\section*{Capítulo \ref{ch:g4:animal-reproduction} --- Como os animais se reproduzem}

\begin{solution}{exo:g4:animal-reproduction:1}
A junção de um \omterm{def:g4:animal-reproduction:sexual}{espermatozoide} com uma \omterm{def:g4:animal-reproduction:sexual}{célula reprodutora feminina},
que dá início a um novo filhote. A \omterm{def:g4:animal-reproduction:sexual}{célula reprodutora feminina} vem
da \omterm{def:g4:animal-reproduction:sexual}{fêmea}, e o \omterm{def:g4:animal-reproduction:sexual}{espermatozoide} vem do \omterm{def:g4:animal-reproduction:sexual}{macho}.
\end{solution}

\begin{solution}{exo:g4:animal-reproduction:2}
\omterm{def:g4:animal-reproduction:ovivivi}{Ovíparo}: a \omterm{def:g4:animal-reproduction:sexual}{fêmea} põe \omterm{def:g2:animals-grow-up:born}{ovos} e o filhote termina de crescer no \omterm{def:g2:animals-grow-up:born}{ovo},
fora do corpo dela --- a galinha, o sapo (também a truta, os
lagartos, os caracóis, os \omterm{prop:g3:sorting-living-things:groups}{insetos}). \omterm{def:g4:animal-reproduction:ovivivi}{Vivíparo}: o filhote cresce
dentro da mãe, alimentado por ela, e nasce de barriga --- a égua, a
baleia (quase todos os \omterm{prop:g3:sorting-living-things:groups}{mamíferos}).
\end{solution}

\begin{solution}{exo:g4:animal-reproduction:3}
Os \omterm{def:g4:animal-reproduction:sexual}{machos} cantaram para chamar as \omterm{def:g4:animal-reproduction:sexual}{fêmeas} até a lagoa; cada casal
soltou as células dele na água, onde os \omterm{def:g4:animal-reproduction:sexual}{espermatozoides} encontraram
os óvulos --- a \omterm{def:g4:animal-reproduction:sexual}{fecundação}; os \omterm{def:g2:animals-grow-up:born}{ovos} fecundados dentro da geleia
viraram os \omterm{ex:g2:animals-grow-up:frog}{girinos} de junho.
\end{solution}

\begin{solution}{exo:g4:animal-reproduction:4}
A da truta acontece na água, fora da \omterm{def:g4:animal-reproduction:sexual}{fêmea} --- óvulos e
\omterm{def:g4:animal-reproduction:sexual}{espermatozoides} são soltos juntos. A da galinha acontece dentro do
corpo dela, durante o \omterm{prop:g4:animal-reproduction:where}{acasalamento}. Em terra as células nuas
secariam; na água elas podem se encontrar em segurança do lado de
fora.
\end{solution}

\begin{solution}{exo:g4:animal-reproduction:5}
O pintinho cresce num \omterm{def:g2:animals-grow-up:born}{ovo} com casca, fora da galinha, vivendo da
gema embalada junto com ele. O potrinho cresce onze meses dentro da
égua, alimentado sem parar pelo corpo dela, e nasce de barriga.
\end{solution}

\begin{solution}{exo:g4:animal-reproduction:6}
Quanto mais filhotes uma espécie produz, menos cuidado cada um
recebe; quanto menos filhotes, mais cuidado. Extremos: os milhões de
\omterm{def:g2:animals-grow-up:born}{ovos} sem cuidado do bacalhau e o filhote único da elefanta,
protegido por uma década.
\end{solution}

\begin{solution}{exo:g4:animal-reproduction:7}
Porque, para dar certo, basta que, em média, sobrevivam filhotes
suficientes para substituir os pais. De milhões de \omterm{def:g2:animals-grow-up:born}{ovos}, alguns
sempre escapam dos perigos do mar --- o número enorme \emph{é} a
proteção.
\end{solution}

\begin{solution}{exo:g4:animal-reproduction:8}
Poucos filhotes (cinco é pouco perto de milhares), crescidos dentro
da mãe, amamentados e protegidos: a gata está claramente do lado dos
poucos e bem cuidados --- como os \omterm{prop:g3:sorting-living-things:groups}{mamíferos} em geral.
\end{solution}

\begin{solution}{exo:g4:animal-reproduction:9}
Que ele recebe alguma coisa de \emph{cada} um dos pais --- o filhote
não é uma cópia só da mãe. Os dois pais contribuem para aquilo que o
potrinho vai ser.
\end{solution}

\begin{solution}{exo:g4:animal-reproduction:10}
Uma desova pequena --- algumas dezenas ou centenas de \omterm{def:g2:animals-grow-up:born}{ovos}, e não
milhões. Cuidado e número se compensam: um pai que protege, abana e
defende consegue levar uma desova pequena através de perigos que
exigiriam milhões de \omterm{def:g2:animals-grow-up:born}{ovos} sem guarda.
\end{solution}

\begin{solution}{exo:g4:animal-reproduction:11}
Ovíparas --- e sem cuidado: a desova enterrada é deixada ao calor da
areia, e os filhotes correm sozinhos até o mar (o número grande
deles é a estratégia). Ela precisa pôr os \omterm{def:g2:animals-grow-up:born}{ovos} em terra porque os
\omterm{def:g2:animals-grow-up:born}{ovos} de réptil têm casca e respiram ar --- postos no mar, eles se
afogariam. Quem impõe a regra são as características da classe dela,
e não os hábitos dela.
\end{solution}
